\documentclass{article}
\usepackage{cite}
\begin{document}

\section{Introduction}

Offshore wind turbines are set to become the most competitive form of wind power generation, which have the advantages of saving land resources and having stronger and steadier wind resources compared with onshore wind turbines \cite{Zhao2025}. Owing to the limited resources of the offshore space, the offshore wind turbine continues to shift from the shallow to deep sea locations. More and more newly-built offshore wind farms with large capacities above 1000 MW will be located further than 100 km offshore \cite{Zhang2025}.

However, the increasing deployment of floating offshore wind turbines in deep-sea environments introduces significant structural challenges. The complex wind-wave coupled loads acting on floating platforms can cause excessive tower and blade vibrations, leading to fatigue damage and reduced operational lifespan \cite{Chen2023, Li2024}. If these structural loads are not adequately mitigated through advanced control strategies, the levelized cost of energy for floating offshore wind will remain prohibitively high.

\subsection{Load Mitigation Approaches}

To address the structural loading problem, various control strategies have been extensively investigated. Active pitch control remains the most widely studied approach for above-rated load reduction \cite{Bossanyi2003, Larsen2007}. Individual pitch control has been developed to reduce asymmetric rotor loads caused by wind shear and tower shadow effects \cite{Bossanyi2005}. More recently, several studies have explored model predictive control for multi-objective load mitigation \cite{Schlipf2013, Kim2019, Abbas2022}.

However, these model-based approaches require accurate system models that are difficult to obtain for floating offshore wind turbines due to the complex nonlinear hydrodynamic interactions. The model uncertainty problem is particularly severe for fault conditions, where actuator degradation fundamentally changes the system dynamics.

\end{document}
